% from ega2-1-fr.tex
\oldpage[II]{115}

\begin{enumerate}
\item[a)] Le point $x$ est isolé dans sa fibre $f^{-1}(f(x))$.
\item[b)] L'anneau $\mathcal{O}_x$ est un $\mathcal{O}_{f(x)}$-module
quasi-fini (\hyperref[0.7.4.1-fr]{0, 7.4.1}).
\end{enumerate}
\end{proposition}

La question étant évidemment locale sur $X$ et sur $Y$, on peut supposer
$X=\operatorname{Spec}(A)$ et $Y=\operatorname{Spec}(B)$ affines, $A$ étant
une $B$-algèbre de type fini (\hyperref[I.6.3.3-fr]{I, 6.3.3}). En outre, on
peut remplacer $X$ par $X\times_Y\operatorname{Spec}(\mathcal{O}_{f(x)})$
sans changer la fibre $f^{-1}(f(x))$ ni l'anneau local $\mathcal{O}_x$
(\hyperref[I.3.6.5-fr]{I, 3.6.5})~; on peut donc supposer que $B$ est un
anneau local (égal à $\mathcal{O}_{f(x)}$)~; si $\mathfrak{n}$ est l'idéal
maximal de $B$, $f^{-1}(f(x))$ est alors un schéma affine d'anneau
$A/\mathfrak{n}A$, de type fini sur $k(f(x))=B/\mathfrak{n}$
(\hyperref[I.6.4.11-fr]{I, 6.4.11}). Cela étant, si a) est vérifiée, on peut
en outre supposer que $f^{-1}(f(x))$ est réduite au point $x$~; donc
$A/\mathfrak{n}A$ est de rang fini sur $B/\mathfrak{n}$
(\hyperref[I.6.4.4-fr]{I, 6.4.4}), autrement dit $A$ est un $B$-module
quasi-fini. Inversement, s'il en est ainsi, $f^{-1}(f(x))$ est un schéma
affine artinien, donc discret (\hyperref[I.6.4.4-fr]{I, 6.4.4})~; par suite
$x$ est isolé dans sa fibre, et cela montre que b) entraîne a).

\begin{corollary}[6.2.2]
\label{II.6.2.2-fr}
Soit $f:X\to Y$ un morphisme de type fini. Les conditions suivantes sont
équivalentes~:
\begin{enumerate}
\item[a)] Tout point $x\in X$ est isolé dans sa fibre $f^{-1}(f(x))$
(autrement dit le sous-espace $f^{-1}(f(x))$ est discret).
\item[b)] Pour tout $x\in X$, le préschéma $f^{-1}(f(x))$ est un
$k(f(x))$-préschéma fini.
\item[c)] Pour tout $x\in X$, l'anneau $\mathcal{O}_x$ est un
$\mathcal{O}_{f(x)}$-module quasi-fini.
\end{enumerate}
\end{corollary}

L'équivalence de a) et c) résulte de
(\hyperref[II.6.2.1-fr]{6.2.1}). D'autre part, comme $f^{-1}(f(x))$ est un
$k(f(x))$-préschéma algébrique (\hyperref[I.6.4.11-fr]{I, 6.4.11}),
l'équivalence de a) et b) résulte de (\hyperref[I.6.4.4-fr]{I, 6.4.4}).

\begin{definition}[6.2.3]
\label{II.6.2.3-fr}
Lorsque $f:X\to Y$ est un morphisme de type fini vérifiant les conditions
équivalentes de (\hyperref[II.6.2.2-fr]{6.2.2}), on dit que $f$ est
\emph{quasi-fini}, ou que $X$ est quasi-fini sur $Y$.
\end{definition}

Il est clair que tout morphisme fini est quasi-fini
(\hyperref[II.6.1.8-fr]{6.1.8}).

\begin{proposition}[6.2.4]
\label{II.6.2.4-fr}
\begin{enumerate}
\item[(i)] Une immersion $X\to Y$, qui est fermée, ou telle que $X$ soit
noethérien, est un morphisme quasi-fini.
\item[(ii)] Si $f:X\to Y$ et $g:Y\to Z$ sont des morphismes quasi-finis,
$g\circ f$ est quasi-fini.
\item[(iii)] Si $X$ et $Y$ sont des $S$-préschémas, $f:X\to Y$ un
$S$-morphisme quasi-fini, alors $f_{(S')}:X_{(S')}\to Y_{(S')}$ est quasi-fini
pour toute extension $g:S'\to S$ de la base.
\item[(iv)] Si $f:X\to Y$ et $g:X'\to Y'$ sont deux $S$-morphismes
quasi-finis,
\[
f\times_S g:X\times_S Y\to X'\times_S Y'
\]
est quasi-fini.
\item[(v)] Soient $f:X\to Y$ et $g:Y\to Z$ deux morphismes tels que $g\circ
f$ soit quasi-fini~; alors, si en outre $g$ est séparé, ou $X$ noethérien, ou
$X\times_ZY$ localement noethérien, $f$ est quasi-fini.
\item[(vi)] Si $f$ est quasi-fini, il en est de même de $f_{\mathrm{red}}$.
\end{enumerate}
\end{proposition}

Si $f:X\to Y$ est une immersion, toute fibre est réduite à un point, et
l'assertion (i) résulte donc de (\hyperref[I.6.3.4-fr]{I, 6.3.4 (i) et
6.3.5}). Pour prouver (ii), on remarque d'abord que $h=g\circ f$ est de type
fini (\hyperref[I.6.3.4-fr]{I, 6.3.4 (ii)})~; en outre, si $z=h(x)$ et
$y=f(x)$, $y$ est isolé dans $g^{-1}(z)$, donc il existe un voisinage ouvert
$V$ de $y$ dans $Y$ ne rencontrant aucun point de $g^{-1}(z)$ distinct de
$y$~; $f^{-1}(V)$ est donc un voisinage ouvert de $x$ ne rencontrant aucun

